% SOURCE_STATEMENT: "Here $h_m$ is the $m$-th Hurwitz class number, i.e. the number of equivalence classes of positive definite binary quadratic forms of discriminant $-m$ with the class containing $x^2+y^2$ weighted by 1/2 and the class containing $x^2 + xy + y^2$ weighted by 1/3. Moreover by convention $h_0=-1/12.$\n\nFor more information in regards to the connection between physics and number theory, see the following work of\n\n Kathrin Bringmann and Ben Kane\nhttp://arxiv.org/pdf/1305.0112v1.pdf\n\nKatrin Bringmann and Sameer Murthy\nhttp://arxiv.org/pdf/1208.3476v2.pdf\n\nKatrin Bringmann and Jan Manschot\nhttp://arxiv.org/pdf/1304.7208v1.pdf\n\nDo the coefficients of $h_m$ have enumerative significance? Can we find $h_m$ for $m \\geq 9$?"
% Preparation and final self-audit: CONTRIBUTING.md prompts C and D; source/audit artifacts retained in the local physics_statements_2026_10_11 batch.
% !TeX program = lualatex
\documentclass[11pt]{article}
\usepackage[margin=25mm]{geometry}
\usepackage{fontspec}
\setmainfont{Latin Modern Roman}
\usepackage{amsmath,amssymb,mathtools}
\usepackage{unicode-math}
\setmathfont{Latin Modern Math}
\usepackage{xurl,hyperref}
\hypersetup{colorlinks=true,urlcolor=blue}
\setcounter{secnumdepth}{0}
\setlength{\parindent}{0pt}
\setlength{\parskip}{5pt}
\setlength{\emergencystretch}{3em}
\newcommand{\sref}[2]{\hyperlink{source:#1:#2}{[S#2]}}
\newcommand{\cataloguescope}[1]{\paragraph{Recorded target.}#1}
\newcommand{\R}{\mathbb R}
\newcommand{\C}{\mathbb C}
\newcommand{\N}{\mathbb N}
\newcommand{\Z}{\mathbb Z}
\newcommand{\Q}{\mathbb Q}
\newcommand{\F}{\mathbb F}
\newcommand{\eps}{\varepsilon}
\begin{document}
% UnsolvedMath ID 20002511; source catalogue code AIM-PHYSICS-0019
\setcounter{section}{20002510}
\section{Hurwitz class numbers, CM enumeration, and local conductor layers}\label{20002511}
{\small UnsolvedMath ID 20002511 · Mathematical Physics}
\subsection{Definitions and mathematical statement}

\paragraph{Shared notation.}$\mathbb N=\{1,2,\ldots\}$; the standard number fields have their usual meanings. All additional conventions are stated locally.


\paragraph{Definitions and notation.}A positive-definite integral binary quadratic form is $ax^2+bxy+cy^2$, with discriminant $b^2-4ac<0$, considered up to $SL_2(\mathbb Z)$ equivalence. Hurwitz class numbers count such classes with automorphism weights, including $1/2$ and $1/3$ for the extra-symmetry classes in the source’s convention. The convention $h_0=-1/12$ is separate from the positive-index class-number counting convention.

\paragraph{Catalogue wording retained for traceability.}
With the source’s Hurwitz class-number convention, including $h_0=-1/12$, determine the enumerative significance of the sequence $h_m$ and how to find $h_m$ for $m\ge9$. The source points to the Bringmann--Kane, Bringmann--Murthy and Bringmann--Manschot works for connections between number theory and physics.
\paragraph{Scope and source qualification.}The source says coefficients of h_m even though h_m already denotes coefficients; this is treated as asking about the sequence. Ordinary finite enumeration is not in itself an unresolved obstruction. The intended additional physical enumerative interpretation remains broad; cited papers were not independently reviewed for current status. The catalogue title/status was paired with a machine-generated research attempt; that attempt is not peer-reviewed evidence of settlement. The statement covers the original source questions rather than only the title’s subcase.
\subsection{Short English statement}
Explain the enumerative meaning of Hurwitz class numbers and compute the values requested in the source.
\subsection{Sources}
\begin{itemize}
\item[\hypertarget{source:20002511:1}{S1}] ULAM.AI and the UnsolvedMath Contributors, UnsolvedMath, record 20002511 (AIM-PHYSICS-0019). Catalogue curation under CC BY 4.0; definitions and English exposition added locally.
\url{https://huggingface.co/datasets/ulamai/UnsolvedMath}.
\item[\hypertarget{source:20002511:2}{S2}] American Institute of Mathematics, Gromov-Witten invariants and number theory, source item 4.2.
\url{http://aimpl.org/gromwitnumthry/4/}.
\item[\hypertarget{source:20002511:3}{S3}] Bringmann and Kane, work linked by the original source; not independently reviewed in this batch.
\url{https://arxiv.org/abs/1305.0112}.
\item[\hypertarget{source:20002511:4}{S4}] Bringmann and Murthy, work linked by the original source; not independently reviewed in this batch.
\url{https://arxiv.org/abs/1208.3476}.
\item[\hypertarget{source:20002511:5}{S5}] Bringmann and Manschot, work linked by the original source; not independently reviewed in this batch.
\url{https://arxiv.org/abs/1304.7208}.
\end{itemize}
\label{end:20002511}
\end{document}
